\documentclass[10pt,a4paper]{amsart}
\usepackage[margin=19mm]{geometry}
\usepackage{amsmath,amssymb,mathtools}
\usepackage[scr=rsfs,cal=euler]{mathalfa}
\usepackage{xfrac}
\usepackage[unicode,hidelinks,bookmarks=false]{hyperref}
\newcommand{\ie}{i.e.\ }
\newcommand{\cf}{cf.\ }
\newcommand{\resp}{respectively\ }
\newcommand{\Subsection}{\subsection}
\providecommand{\nobreakdash}{\nobreak\hbox{-}}
\setlength{\emergencystretch}{2em}
\multlinegap0pt
\usepackage{bm}
\usepackage{bbm}
\usepackage{dsfont}
\usepackage{xspace}
\usepackage{cancel}
\usepackage{mathtools}
\usepackage{amsthm}
\makeatletter
\@ifundefined{thm}{\newtheorem{thm}{Thm}}{}
\@ifundefined{prop}{\newtheorem{prop}[thm]{Proposition}}{}
\makeatother
\makeatletter
\newcommand{\R}{{\mathbb R}}
\newcommand{\Z}{{\mathbb Z}}
\newcommand{\length}{\operatorname{length}}
\expandafter\ifx\csname proofmthm\endcsname\relax
\newenvironment{proofmthm}{\begin{proof}[Proof of the Main Theorem]}
                   {\end{proof}}
\fi
\expandafter\ifx\csname say\endcsname\relax
\newenvironment{say}{}{}
\fi
\expandafter\ifx\csname sketch\endcsname\relax
\newenvironment{sketch}{\begin{proof}[Sketch of proof]}{\end{proof}}
\fi
\makeatother
\title[An explicit bound for complements on toric varieties]{An explicit bound for complements on toric varieties}
\author{Florin Ambro}
\date{Source: arXiv:2506.18509v2 (CC BY 4.0)}
\begin{document}
\maketitle

\noindent\emph{Source and licence.} An explicit bound for complements on toric varieties. Version arXiv:2506.18509v2 (CC BY 4.0), distributed under the Creative Commons Attribution 4.0 International licence, as recorded in the arXiv licence metadata for this exact version and shown on the abstract page https://arxiv.org/abs/2506.18509v2. Full rights notice: licenses/arxiv-2506.18509-CC-BY-4.0.txt.\par\smallskip

\noindent\textit{Editorial scope.} Counted unit: the Prop whose source label is \texttt{lfb} in the article (source0.tex lines 404--413, source label \texttt{lfb}), delivered together with the article's own complete original proof of it (source0.tex lines 415--443, one \texttt{proof} environment of 29 lines). No mathematics is written, repaired or re-derived: statement and proof are the article's own text line for line. Required context retained uncounted: none. Every definition, display and label of those lines is retained unchanged. Omitted from this package and not redistributed: 1--403, 414--414, 444--678. Nothing inside a retained excerpt is omitted or altered: every retained body is delivered exactly as the article's lines are written, so no omission receipt of any kind accompanies this package. Citations: the retained excerpts carry no citation command at all, so no bibliography entry of the article is transcribed and no reference is redistributed. Reference adaptation: none was needed and none was made; every reference inside the delivered unit resolves inside it, so no unresolved reference or citation is produced. Printed slips kept literal: no printed word, symbol, hypothesis or number of the counted statement or of its proof was altered. Layout and class: the article's own class is \texttt{amsart} with packages including amsthm, amsfonts, amssymb, amscd, euscript, xy; the wrapper is the lane's pinned \texttt{amsart} editorial wrapper at 10pt on A4, which restates the theorem-like environments the retained text needs and adds the article's own macro definitions that the retained text needs (\texttt{\textbackslash R}, \texttt{\textbackslash Z}, \texttt{\textbackslash length}), with extra wrapper packages (bm, bbm, dsfont, xspace, cancel, mathtools). The delivered numbering is the unit's own. Assets: no retained span contains a figure environment, a graphics-inclusion command, a PDF-page inclusion, a table or a TikZ picture; no image file is copied into this package, the package declares no asset dependency and no image file is redistributed with it. Licence: version arXiv:2506.18509v2 of this article is distributed under the Creative Commons Attribution 4.0 International licence, as recorded in the arXiv licence metadata for this exact version and shown on the abstract page https://arxiv.org/abs/2506.18509v2. A full rights notice is delivered at licenses/arxiv-2506.18509-CC-BY-4.0.txt. Characters: every retained excerpt is pure ASCII, so the delivered engine, XeLaTeX with its default Latin Modern text font, reports no missing character.
\begin{prop}\label{lfb}
	Let $\varphi_0\colon \Lambda_0\to \Z$ be a surjective homomorphism.
	Suppose the $\varphi_0$-image of a translation of $C \cap\varphi^{-1}(p)$ which is contained in $\varphi^\perp$
	has finite positive length $w_0$. Then there exists
	a surjective homomorphism $\varphi'\colon \Lambda\to \Z$ such that $\varphi'|_{\Lambda_0}=\varphi_0$
	and 
	$
	\length \varphi'(C)<w(\varphi)+\frac{w_0}{\gamma}.
	$
\end{prop}

\begin{proof}
	Let $\varphi'\colon \Lambda\to \Z$ be an arbitrary extension of $\varphi_0$. Then 
	$\Phi=(\varphi,\varphi')\colon \Lambda\to \Z^2$ is surjective, $\Phi(C)$ is
	a convex body in $\R^2$ such that its projection onto the $x_1$-axis is the interval $[p_1,p_2]$, and its intersection with the line $\{x_1=p\}$ is an interval of length $w_0$. Through each endpoint of the interval $\Phi(C)\cap \{x_1=p\}$, draw a supporting hyperplane to $\Phi(C)$. Together with the inequalities $p_1 \le x_1 
	\le p_2$, these form a quadrilateral or triangle $Q$, which contains $\Phi(C)$. We minimize and maximize the $x_2$-coordinate among the vertices of $Q$. There are two cases.
	
	1) Case both min and max are attained on the same vertical line defining $Q$. Say $(p_2,y)$ and $(p_2,y')$ are such that $[y,y']$ is the projection of $Q$ onto the $x_2$-axis. Then $y'-y=\length  \varphi'(Q)=:d$. Keeping the same intersection $Q\cap \{x_1=p\}$, we slide one of the vertices of $Q\cap\{x_1=p_1\}$ towards the other until $Q$ becomes a triangle $Q'$ with a vertex on the line $x_1=p_1$ and base a segment on the line $x_1=p_2$ of length $d'\ge d$. By similarity, we compute
	$$
	\frac{d'}{w_0}=\frac{p_2-p_1}{p-p_1}=\frac{1}{\lambda_2}.
	$$
	Therefore $d\le d'=\frac{w_0}{\lambda_2}\le \frac{w_0}{\gamma}$. So in this case, $\varphi'$ is the desired extension, satisfying the stronger inequality 
	$
	\length \varphi'(C) \le \frac{w_0}{\gamma}.
	$
	
	2) Case min and max are attained on different vertical lines defining $Q$. Say $(p_1,y)$ and $(p_2,y')$ are vertices of $Q$ such that $[y,y']$ is the projection of $Q$ onto the $x_2$-axis. Let $(p,y'')$ be the vertex of $Q\cap \{x_1=p\}$ with maximal $x_2$-value.	
	Let $\delta=y'-y''$. Keeping the vertex $(p_2,y')$ of $Q$ fixed, we rotate the non-vertical supporting hyperplane of $Q$ until it becomes $\{x_2=y'\}$. The new body $Q'$ has again $[y,y']$ as projection onto the $x_2$-axis, $(p_1,y),(p_1,y')$ are among its vertices, and its intersection with $\{x_1=p\}$ has length $w_0+\delta$. By the argument of case 1), we have 
	$$
	\frac{\length \varphi'(C)}{w_0+\delta}\le \frac{1}{\lambda_1}.
	$$ 	
	For $z\in \Z$, we have an automorphism of $\Z^2$ defined by $x'_1=x_1,x'_2=zx_1+x_2$. It acts on each line $\{x_1=a\}$ as the translation $y\mapsto y+za$.
	After this automorphism, $\delta$ transforms into $\delta+z(p_2-p)$.
	There exists a unique $z\in \Z$ such that $\delta+z(p_2-p)\in [0,p_2-p)$, namely 
	$z=\lceil \frac{-\delta}{p_2-p}\rceil$. Therefore, after changing the second coordinate
	(i.e. changing the lifting of $\varphi_0$ from $\varphi'$ to $\varphi'+z\varphi$), we may suppose $\delta<p_2-p=\lambda_1(p_2-p_1)$. Then
	$
	\length \varphi'(C) <\frac{w_0}{\lambda_1}+p_2-p_1\le \frac{w_0}{\gamma}+w(\varphi).
	$
\end{proof}

\end{document}
